\documentclass[12pt,a4paper]{article}

\usepackage[utf8]{inputenc}
\usepackage[T1]{fontenc}
\usepackage[french]{babel}
\usepackage{geometry}
\usepackage{booktabs}
\usepackage{hyperref}
\usepackage{tabularx}
\usepackage{enumitem}

\geometry{margin=2.5cm}

\hypersetup{
    colorlinks=true,
    linkcolor=blue!60!black,
    urlcolor=blue!60!black
}

\title{Synthèse par membre \\ \Large Projet Bloc IA --- HumanForYou}
\author{
    Maxence Chartier \and
    Amine Haouri \and
    Anisse Imerzoukene \and
    Mourad Messaoudi
}
\date{\today}

\begin{document}

\maketitle
\tableofcontents
\newpage

% ============================================================
\section{Répartition des tâches}

\begin{table}[h!]
\centering
\caption{Répartition des responsabilités par membre}
\small
\begin{tabularx}{\textwidth}{l X}
\toprule
\textbf{Membre} & \textbf{Responsabilités principales} \\
\midrule
Maxence Chartier & Cadre réglementaire européen, prétraitement des données \\
Amine Haouri & Cadre réglementaire indien et comparatif, exploration des données \\
Anisse Imerzoukene & Description des données, fusion des datasets, modélisation \\
Mourad Messaoudi & Enjeux éthiques appliqués au projet, résultats et recommandations \\
\bottomrule
\end{tabularx}
\end{table}

% ============================================================
\section{Maxence Chartier}

\subsection{Synthèse des sources étudiées}

\paragraph{RGPD (CNIL).}
Le Règlement Général sur la Protection des Données impose six principes pour tout traitement de données personnelles : licéité, limitation des finalités, minimisation, exactitude, limitation de conservation, et sécurité. L'article~22 interdit les décisions fondées uniquement sur un traitement automatisé ayant des effets significatifs sur une personne, sauf exceptions encadrées.

\paragraph{AI Act (Règlement UE 2024/1689).}
L'AI Act classe les systèmes d'IA par niveaux de risque. Les systèmes RH (recrutement, évaluation, prédiction d'attrition) sont classés \textbf{haut risque}. Obligations : gestion des risques, gouvernance des données, documentation, contrôle humain, analyse d'impact.

\paragraph{Guide de sécurité des données personnelles (CNIL, 2023).}
Guide pratique pour la mise en conformité technique : mesures de sécurité (chiffrement, contrôle d'accès, journalisation), gestion des violations, sensibilisation des équipes.

\subsection{Contribution au projet}

Rédaction de la partie réglementaire européenne (RGPD, AI Act, principes du HLEG). Participation au prétraitement : gestion des valeurs manquantes, suppression des variables constantes (\texttt{EmployeeCount}, \texttt{Over18}, \texttt{StandardHours}), encodage des variables catégorielles.

% ============================================================
\section{Amine Haouri}

\subsection{Synthèse des sources étudiées}

\paragraph{Digital Personal Data Protection Act (DPDPA, 2023).}
Première loi cadre indienne de protection des données numériques. Deux bases juridiques : consentement et utilisations légitimes. Le traitement des données employés est autorisé sans consentement pour les finalités d'emploi (Section~7(i)). Sanctions jusqu'à 2,5~milliards INR.

\paragraph{India AI Governance Guidelines (2025).}
Approche \textit{light-touch} : pas de loi dédiée à l'IA, mais des lignes directrices non contraignantes. Quatre principes : confiance, approche centrée sur l'humain, innovation, équité. Architecture institutionnelle en cours de mise en place (AIGG, TPEC, AISI, NCAIC).

\paragraph{Analyse comparative UE/Inde.}
L'UE impose des obligations contraignantes avec classification par risque ; l'Inde privilégie la flexibilité et l'innovation. Convergence sur les principes (transparence, équité, sécurité) ; divergence sur le caractère contraignant et les sanctions.

\subsection{Contribution au projet}

Rédaction de la partie réglementaire indienne et de l'analyse comparative. Exploration des données : statistiques descriptives, distributions, identification des valeurs aberrantes, analyse du taux d'attrition par catégorie.

% ============================================================
\section{Anisse Imerzoukene}

\subsection{Synthèse des sources étudiées}

\paragraph{Rapport Villani (2018).}
Rapport stratégique de 235~pages définissant la stratégie française en IA autour de six axes : politique économique des données, recherche, emploi, écologie, éthique et inclusivité. Précurseur de l'approche européenne actuelle.

\paragraph{Lignes directrices du HLEG (2019) et Considérant~27 de l'AI Act.}
Sept principes éthiques pour une IA digne de confiance : intervention humaine, robustesse, vie privée, transparence, non-discrimination, bien-être sociétal, responsabilité. Ces principes, initialement non contraignants, sont désormais intégrés au cadre de l'AI Act.

\paragraph{Guide pratique IA éthique (Hub France IA, 2024).}
Guide opérationnel pour la mise en \oe uvre d'une démarche d'IA éthique en entreprise : méthodologie, outils d'évaluation, bonnes pratiques.

\subsection{Contribution au projet}

Description détaillée des quatre sources de données. Développement du script de fusion (\texttt{build\_dataset.py}) et production du \texttt{dataset\_final.csv}. Choix et entraînement des modèles de prédiction, validation croisée, recherche d'hyperparamètres.

% ============================================================
\section{Mourad Messaoudi}

\subsection{Synthèse des sources étudiées}

\paragraph{Rapport CNIL sur l'éthique des algorithmes (2017).}
Deux principes fondateurs : loyauté (les algorithmes au service des citoyens) et vigilance/réflexivité (questionnement continu). Six recommandations : formation à l'éthique, compréhensibilité, liberté humaine, plateforme d'audit, recherche, fonction éthique en entreprise.

\paragraph{CEP -- Analyse des lignes directrices éthiques (COM 2019/168).}
Analyse critique de la communication de la Commission européenne sur les lignes directrices éthiques. Mise en perspective des enjeux de mise en \oe uvre pratique des principes éthiques dans le contexte industriel.

\paragraph{Lignes directrices et recommandations de la CNIL.}
Cadre méthodologique pour les analyses d'impact (AIPD), recommandations sur les cookies et traceurs, boîte à outils RGPD pour les professionnels.

\subsection{Contribution au projet}

Analyse des enjeux éthiques spécifiques au projet HumanForYou : risques de biais (genre, âge, statut marital), transparence et explicabilité du modèle, supervision humaine des décisions. Interprétation des résultats et formulation des recommandations stratégiques à destination de la direction.

% ============================================================
\section{Synthèse collective}

\begin{table}[h!]
\centering
\caption{Sources étudiées par membre}
\small
\begin{tabularx}{\textwidth}{l X}
\toprule
\textbf{Membre} & \textbf{Sources principales} \\
\midrule
Maxence & RGPD (CNIL), AI Act, Guide sécurité CNIL 2023 \\
Amine & DPDPA 2023, India AI Guidelines 2025, Comparatif UE/Inde \\
Anisse & Rapport Villani, HLEG / Considérant 27, Hub France IA \\
Mourad & Rapport CNIL 2017, CEP COM 2019/168, Lignes directrices CNIL \\
\bottomrule
\end{tabularx}
\end{table}

\paragraph{Travail commun.}
L'ensemble du groupe a participé aux phases d'exploration, de modélisation et de rédaction finale. La répartition ci-dessus reflète les responsabilités principales, chaque membre ayant contribué à la relecture et à l'amélioration de l'ensemble du document.

\end{document}
